\documentclass[11pt]{article}
\usepackage[margin=1.05in]{geometry}
\usepackage{amsmath,amssymb,amsthm,mathtools}
\usepackage{booktabs}
\usepackage{hyperref}

\newtheorem{theorem}{Theorem}
\newtheorem{proposition}[theorem]{Proposition}
\newtheorem{lemma}[theorem]{Lemma}
\newtheorem{corollary}[theorem]{Corollary}
\theoremstyle{definition}
\newtheorem{definition}[theorem]{Definition}
\newtheorem{remark}[theorem]{Remark}
\newtheorem{example}[theorem]{Example}

\newcommand{\bb}[1]{[\![#1]\!]}
\newcommand{\lam}{\lambda}
\newcommand{\Qp}{Q'}
\newcommand{\Lam}{\Lambda}

\title{What $M^{(L)}$ is:\\ the length-filtration matrix of $\langle p_\rho,h_\nu\rangle$}
\author{Clio}
\date{8 October 2026}

\begin{document}
\maketitle

\begin{abstract}
Let $M$ be the matrix $M_{\rho\nu}=\langle p_\rho,h_\nu\rangle=[m_\nu]\,p_\rho$ on partitions of
$n$, and let $M^{(L)}$ be its leading block on $\{\rho:\ell(\rho)\le L\}$. In earlier work I proved
that $M$ is lower block-triangular for length, so that every $M^{(L)}$ is invertible; the question
left open was what $M^{(L)}$ actually \emph{is}. The answer is that the filtration is much finer
than block-triangular: \textbf{each diagonal block is a diagonal matrix}, with entry
$\prod_i m_i(\rho)!$ at $\rho$. Consequently
\[
  \det M^{(L)}\;=\;\prod_{\substack{\rho\vdash n\\ \ell(\rho)\le L}}\ \prod_{i\ge 1} m_i(\rho)!
  \;=\;\prod_{\substack{\rho\vdash n\\ \ell(\rho)\le L}}\frac{z_\rho}{\rho_1\cdots\rho_{\ell(\rho)}},
\]
which is $1$ exactly when $L=1$, or $n=1$, or $L=2$ with $n$ odd. The inverse is the Möbius function
of the partition lattice, it respects the same filtration, and the row of $M^{-1}$ indexed by $\nu$
has denominator exactly $\prod_k m_k(\nu)!$ --- so the $h$-coefficients are recovered from the
$p$-coefficients by dividing by automorphism orders and nothing else. I give $(M^{(3)})^{-1}$
explicitly and deduce the $L=3$ analogue of Theorem~B: a three-part class of a Green polynomial sees
at most \emph{five} of the $p(n)$ coefficients $c_{\lam,\nu}$, and in general a $k$-part class sees
at most $B_k$ of them, the $k$-th Bell number.
\end{abstract}

\section{The object}

Work in $\Lam_n$, the degree-$n$ part of the ring of symmetric functions over $\mathbb{Q}$, with the
Hall inner product, for which $\langle p_\rho,p_\sigma\rangle=\delta_{\rho\sigma}z_\rho$ and
$\langle h_\nu,m_\sigma\rangle=\delta_{\nu\sigma}$. For a partition $\rho$ write $\ell(\rho)$ for
its number of parts, $m_i(\rho)$ for the multiplicity of $i$ among its parts,
$z_\rho=\prod_i i^{m_i(\rho)}m_i(\rho)!$, and
\[
  d_\rho\;\coloneqq\;\prod_{i\ge1} m_i(\rho)!\;=\;\frac{z_\rho}{\rho_1\rho_2\cdots\rho_{\ell(\rho)}}.
\]
Thus $d_\rho$ is the order of the group of permutations of the \emph{positions} of $\rho$ that fix
the sequence of parts: the automorphism group of $\rho$ viewed as a labelled multiset.

\begin{definition}\label{def:M}
$M=M(n)$ is the $p(n)\times p(n)$ matrix with rows and columns indexed by partitions of $n$ and
\[
  M_{\rho\nu}\;=\;\langle p_\rho,\,h_\nu\rangle .
\]
For $1\le L\le n$, $M^{(L)}$ is the submatrix on
$\{\rho\vdash n:\ell(\rho)\le L\}\times\{\nu\vdash n:\ell(\nu)\le L\}$.
\end{definition}

Two further readings of the same entry will be used throughout. Both are standard; I record them
because the proofs below use all three faces of the object.

\begin{lemma}\label{lem:threefaces}
For all $\rho,\nu\vdash n$,
\[
  \langle p_\rho,h_\nu\rangle
  \;=\;[m_\nu]\,p_\rho
  \;=\;\#\Bigl\{f:[\ell(\rho)]\to[\ell(\nu)]\ \Bigm|\ \textstyle\sum_{i\in f^{-1}(j)}\rho_i=\nu_j
    \ \text{for all }j\Bigr\}.
\]
In particular $M$ is a matrix of non-negative integers, and $M_{\rho\nu}$ is the number of ordered
set partitions $(B_1,\dots,B_{\ell(\nu)})$ of the index set of the parts of $\rho$ with
$\sum_{i\in B_j}\rho_i=\nu_j$.
\end{lemma}

\begin{proof}
The first equality is $\langle h_\nu,m_\sigma\rangle=\delta_{\nu\sigma}$ applied to the
$m$-expansion of $p_\rho$. For the second, $p_\rho=\prod_{i=1}^{\ell(\rho)}\bigl(\sum_v
x_v^{\rho_i}\bigr)$, so expanding the product indexes the terms by the choice, for each $i$, of a
variable $x_{v_i}$; the resulting monomial is $\prod_v x_v^{\,s_v}$ with
$s_v=\sum_{i:v_i=v}\rho_i$. Collecting the terms whose exponent multiset is $\nu$ and specialising
to the monomial $x_1^{\nu_1}\cdots x_{\ell(\nu)}^{\nu_{\ell(\nu)}}$ identifies the count with the
number of maps $f:i\mapsto v_i$ onto the prescribed exponents, i.e.\ with the displayed set.
\end{proof}

\begin{remark}\label{rem:nonempty}
Every part $\nu_j$ is $\ge1$ and every $\rho_i$ is $\ge1$, so the condition
$\sum_{i\in f^{-1}(j)}\rho_i=\nu_j$ already forces $f^{-1}(j)\ne\emptyset$. The surjectivity of $f$
is therefore not an extra hypothesis but a consequence --- a point that matters in
\S\ref{sec:verif}, where the obvious control built on it turns out to be vacuous, and in
Remark~\ref{rem:loadbearing}, where positivity of the parts of $\nu$ is shown to be load-bearing
for Theorem~\ref{thm:structure}.
\end{remark}

\section{The structure theorem}

\begin{theorem}[structure of $M^{(L)}$]\label{thm:structure}
Let $\rho,\nu\vdash n$.
\begin{enumerate}
\item[(i)] If $\ell(\nu)>\ell(\rho)$ then $M_{\rho\nu}=0$.
\item[(ii)] If $\ell(\nu)=\ell(\rho)$ and $\nu\ne\rho$ then $M_{\rho\nu}=0$.
\item[(iii)] $M_{\rho\rho}=d_\rho=\prod_i m_i(\rho)!$.
\end{enumerate}
Consequently, for any linear order on partitions of $n$ refining the order by $\ell$, the matrix $M$
is lower triangular with diagonal $(d_\rho)_\rho$; the same holds for every leading block $M^{(L)}$,
and $M^{(L)}$ is exactly the submatrix of $M$ cut out by $\ell\le L$.
\end{theorem}

\begin{proof}
Fix $\rho,\nu$ and use the third description of Lemma~\ref{lem:threefaces}. By
Remark~\ref{rem:nonempty} every such $f$ is surjective, so $\ell(\nu)\le\ell(\rho)$: this is (i).

Now suppose $\ell(\nu)=\ell(\rho)=k$. A surjection between two finite sets of the same cardinality
$k$ is a bijection. So each $f$ counted is a bijection $[k]\to[k]$, every fibre is a singleton, and
the constraint reads $\rho_i=\nu_{f(i)}$ for all $i$. Hence such an $f$ exists only if
$(\rho_1,\dots,\rho_k)$ and $(\nu_1,\dots,\nu_k)$ agree as multisets; both being weakly decreasing,
that forces $\nu=\rho$, which is (ii). When $\nu=\rho$, the $f$ counted are precisely the bijections
of $[k]$ with $\rho_{i}=\rho_{f(i)}$ for all $i$ --- the automorphisms of $\rho$ as a labelled
multiset --- and there are $\prod_i m_i(\rho)!$ of them. That is (iii).

The final sentence is immediate: (i) and (ii) say $M_{\rho\nu}=0$ unless $\ell(\nu)<\ell(\rho)$ or
$\nu=\rho$, so in any order refining $\ell$ all nonzero off-diagonal entries lie strictly below the
diagonal. The set $\{\ell\le L\}$ is an initial segment of such an order, so $M^{(L)}$ is a leading
principal submatrix.
\end{proof}

\begin{remark}
The whole content is the sentence \emph{a surjection between finite sets of equal cardinality is a
bijection}. The earlier proof of (i) stopped at ``block-triangular''; the step from there to (ii) is
free, and it is (ii) that makes the determinant a product and the inverse computable blockwise. This
is the shape answer the question asked for: the length filtration on $\Lam_n$ does not merely
triangularise the $p$-vs-$h$ pairing, it \emph{diagonalises} it on each graded piece.
\end{remark}

\begin{remark}[positivity of $\nu$ is load-bearing]\label{rem:loadbearing}
Part (ii) fails the moment zero parts are admitted. Extend the column index to weak compositions:
for $\rho=(3,2)$ and $\nu=(5,0)$ both of length $2$ and $\rho\ne\nu$, the count of maps
$[2]\to[2]$ with block sums $(5,0)$ is $1$ (send both parts to block~$1$), not $0$. So (ii) is a
statement about partitions, not about matrices of this general shape. This is a fact about the
statement, not about the proof.
\end{remark}

\section{The determinant}

\begin{corollary}\label{cor:det}
For every $n\ge1$ and $1\le L\le n$,
\[
  \boxed{\ \det M^{(L)}\;=\;\prod_{\substack{\rho\vdash n\\ \ell(\rho)\le L}} \prod_{i\ge1}m_i(\rho)!
  \;=\;\prod_{\substack{\rho\vdash n\\ \ell(\rho)\le L}}\frac{z_\rho}{\rho_1\cdots\rho_{\ell(\rho)}}\ }
\]
In particular $\det M^{(1)}=1$, $M^{(L)}$ is invertible over $\mathbb{Q}$ for every $L$, and
\[
  \det M^{(2)}=1+\bb{n\text{ even}} .
\]
\end{corollary}

\begin{proof}
By Theorem~\ref{thm:structure} the matrix is triangular with diagonal $(d_\rho)$. For $L=1$ the only
index is $\rho=(n)$, with $d_{(n)}=m_n!=1$. For $L=2$ the indices are $(n)$ and the two-part
partitions $(a,b)$, $a\ge b\ge1$; $d_{(a,b)}=1$ unless $a=b$, when it is $2!=2$, and $(a,a)$ occurs
among partitions of $n$ exactly when $n$ is even.
\end{proof}

\begin{corollary}[when the change of basis is unimodular]\label{cor:unimod}
$\det M^{(L)}=1$ if and only if every $\rho\vdash n$ with $\ell(\rho)\le L$ is multiplicity-free,
if and only if
\[
  L=1,\qquad\text{or}\qquad n=1,\qquad\text{or}\qquad (L=2 \text{ and } n \text{ odd}).
\]
\end{corollary}

\begin{proof}
$\det M^{(L)}=1$ iff $d_\rho=1$ for all $\rho$ in range, i.e.\ iff no such $\rho$ has a repeated
part. For $L=1$ or $n=1$ there is nothing to check. For $L=2$: the only partition of $n$ into at most
two parts with a repeat is $(n/2,n/2)$, which exists iff $n$ is even. For $L\ge3$ and $n\ge3$ the
partition $(n-2,1,1)$ --- read as $(1,1,1)$ when $n=3$ --- has at most three parts and a repeated
part; for $L\ge2$ and $n=2$ the partition $(1,1)$ does. So the listed cases are exactly the
unimodular ones.
\end{proof}

\begin{remark}
Taking $L=n$ gives $\det M=\prod_{\rho\vdash n}d_\rho$, the sequence
$1,2,6,96,2880,9953280,\dots$ for $n=1,2,\dots$. This has an independent check that uses none of the
machinery above. Since $M(p,s)=M(p,m)\,M(m,s)$ and the Kostka matrix $M(s,m)$ is unitriangular for
dominance order, $|\det M(m,s)|=1$, so $|\det M|$ must equal the absolute value of the determinant
of the character table of $S_n$. It does, for every $n\le7$ (\S\ref{sec:verif}, D).
\end{remark}

\section{The inverse, and why the filtration survives it}

\begin{theorem}[Möbius form of the inverse]\label{thm:inverse}
Let $\Pi_k$ be the lattice of set partitions of $[k]$. For $\lam\vdash n$ with $\ell(\lam)=k$ and
$\pi\in\Pi_k$ write $\lam^\pi$ for the partition of $n$ whose parts are the block sums
$\bigl(\sum_{i\in B}\lam_i\bigr)_{B\in\pi}$. Then
\[
  p_\lam\;=\;\sum_{\pi\in\Pi_{k}}\widetilde m_{\lam^\pi},
  \qquad\text{and}\qquad
  (M^{-1})_{\lam\nu}\;=\;\frac{1}{d_\lam}\sum_{\substack{\pi\in\Pi_{k}\\ \lam^\pi=\nu}}\ \prod_{B\in\pi}(-1)^{|B|-1}(|B|-1)!,
\]
where $\widetilde m_\mu=d_\mu\,m_\mu$ is the augmented monomial symmetric function.
\end{theorem}

\begin{proof}
Write $\widetilde m_{(a_1,\dots,a_r)}=\sum x_{v_1}^{a_1}\cdots x_{v_r}^{a_r}$, summed over tuples of
\emph{distinct} variables; this depends only on the multiset $\{a_i\}$ and equals $d_\mu m_\mu$ for
$\mu=\mathrm{sort}(a)$, since exactly $d_\mu$ orderings of the exponents give each monomial. Now
$p_\lam=\sum_{v_1,\dots,v_k}x_{v_1}^{\lam_1}\cdots x_{v_k}^{\lam_k}$ over \emph{all} tuples; sorting
the tuples by the set partition $\pi$ of $[k]$ recording which $v_i$ coincide, and summing the
exponents within each block, gives $p_\lam=\sum_{\pi\in\Pi_k}\widetilde m_{\lam^\pi}$ as claimed.

This is a sum over the interval $[\hat0,\hat1]$ of $\Pi_k$ in disguise: define
$F(\pi)=\widetilde m_{\lam^\pi}$ and $G(\pi)=\sum_{\sigma\ge\pi}F(\sigma)$. The displayed identity
applied to each block of $\pi$ separately gives $G(\pi)=\prod_{B\in\pi}p_{\lam|_B}$ --- the
power-sum product over the blocks of $\pi$ --- and $F(\hat0)=\widetilde m_\lam$. Möbius inversion in
$\Pi_k$ therefore yields
\[
  \widetilde m_\lam\;=\;\sum_{\pi\in\Pi_k}\mu(\hat0,\pi)\,\prod_{B\in\pi}p_{\lam|_B}
  \;=\;\sum_{\pi\in\Pi_k}\Bigl(\prod_{B\in\pi}(-1)^{|B|-1}(|B|-1)!\Bigr)\,p_{\lam^\pi},
\]
using $\mu(\hat0,\pi)=\prod_{B\in\pi}(-1)^{|B|-1}(|B|-1)!$ for the partition lattice. Dividing by
$d_\lam$ expresses $m_\lam$ in the $p$-basis. Finally, $M$ is the transition matrix
$p_\lam=\sum_\nu M_{\lam\nu}m_\nu$ (Lemma~\ref{lem:threefaces}), so the matrix expressing $m$ in
terms of $p$ is exactly $M^{-1}$, and reading off the coefficient of $p_\nu$ gives the formula.
\end{proof}

\begin{remark}
Theorem~\ref{thm:inverse} is classical --- it is the partition-lattice form of the $m$-to-$p$
transition, in the spirit of Doubilet's Möbius-function treatment of the transition matrices. I
reprove it here because what I need from it is not the formula but its \emph{support}, and because
the normalisation $1/d_\lam$ is precisely the quantity Theorem~\ref{thm:structure}(iii) produced
from the other side.
\end{remark}

\begin{corollary}[the filtration survives inversion]\label{cor:truncinv}
$(M^{-1})_{\lam\nu}=0$ unless $\ell(\nu)\le\ell(\lam)$, and for every $L$
\[
  \bigl(M^{(L)}\bigr)^{-1}\;=\;\bigl(M^{-1}\bigr)^{(L)} .
\]
\end{corollary}

\begin{proof}
A set partition $\pi\in\Pi_{\ell(\lam)}$ has at most $\ell(\lam)$ blocks, so
$\ell(\lam^\pi)\le\ell(\lam)$ and the sum in Theorem~\ref{thm:inverse} is empty unless
$\ell(\nu)\le\ell(\lam)$. The second statement is the general fact that for a block-triangular
$T=\begin{psmallmatrix}A&0\\ C&D\end{psmallmatrix}$ one has
$T^{-1}=\begin{psmallmatrix}A^{-1}&0\\ -D^{-1}CA^{-1}&D^{-1}\end{psmallmatrix}$, whose leading block
is $A^{-1}$; here $A=M^{(L)}$ because $\{\ell\le L\}$ is an initial segment.
\end{proof}

\begin{corollary}[exact denominators]\label{cor:denom}
Every entry of the row of $M^{-1}$ indexed by $\nu$ lies in $\frac{1}{d_\nu}\mathbb{Z}$, and the
least common denominator of that row is exactly $d_\nu$. Hence if
$Y_\rho=\sum_\nu M_{\rho\nu}c_\nu$, then $c\mapsto Y$ is integral, while
\[
  d_\nu\,c_\nu\ \in\ \mathbb{Z}\text{-span}\bigl\{\,Y_\rho:\ \ell(\rho)\le\ell(\nu)\,\bigr\},
\]
and $d_\nu$ cannot be replaced by any proper divisor.
\end{corollary}

\begin{proof}
Theorem~\ref{thm:inverse} exhibits each entry of row $\nu$ as $1/d_\nu$ times an integer;
Theorem~\ref{thm:structure}(iii) and Corollary~\ref{cor:truncinv} give the diagonal entry
$(M^{-1})_{\nu\nu}=1/d_\nu$ exactly, whose denominator is already $d_\nu$. The support statement is
Corollary~\ref{cor:truncinv}.
\end{proof}

\begin{remark}
Corollary~\ref{cor:denom} is the arithmetic content of the length-filtration duality, and it is sharp
in both directions: passing from the $h$-coefficients to the $p$-class values costs nothing, and
passing back costs exactly the automorphism orders of the classes --- no primes beyond those dividing
$\prod_k m_k(\nu)!$ ever appear. By Corollary~\ref{cor:unimod} the duality is an isomorphism of
\emph{integer} lattices only in the three listed degenerate cases.
\end{remark}

\section{$(M^{(3)})^{-1}$ explicitly}

\begin{theorem}\label{thm:M3inv}
Index rows and columns by $\{\rho\vdash n:\ell(\rho)\le3\}$ and write $d_\rho=\prod_k m_k(\rho)!$.
Then $N=(M^{(3)})^{-1}$ is given by
\[
N_{\rho\nu}=
\begin{cases}
1, & \ell(\rho)=1,\ \nu=(n),\\[2pt]
-\dfrac{1}{d_\rho}, & \ell(\rho)=2,\ \nu=(n),\\[6pt]
\dfrac{2}{d_\rho}, & \ell(\rho)=3,\ \nu=(n),\\[6pt]
-\dfrac{m_{\nu_1}(\rho)+m_{\nu_2}(\rho)}{d_\rho\bigl(1+\bb{\nu_1=\nu_2}\bigr)},
  & \ell(\rho)=3,\ \ell(\nu)=2,\\[8pt]
\dfrac{\delta_{\rho\nu}}{d_\rho}, & \ell(\rho)=\ell(\nu)\in\{2,3\},\\[6pt]
0, & \ell(\nu)>\ell(\rho).
\end{cases}
\]
\end{theorem}

\begin{proof}
Write $M^{(3)}=\begin{psmallmatrix}D_1&0&0\\ A_{21}&D_2&0\\ A_{31}&A_{32}&D_3\end{psmallmatrix}$ in
the length grading, so that by Theorem~\ref{thm:structure} each $D_k$ is the diagonal matrix
$\mathrm{diag}(d_\rho)_{\ell(\rho)=k}$. The off-diagonal blocks are read off
Lemma~\ref{lem:threefaces}:
\begin{itemize}
\item $A_{k1}$: for $\nu=(n)$ there is exactly one map $[\,\ell(\rho)]\to[1]$ and its block sum is
$n$, so every entry of the first column is $1$.
\item $A_{32}$: let $\ell(\rho)=3$ and $\ell(\nu)=2$. A surjection $[3]\to[2]$ has one singleton
fibre and one fibre of size $2$; the two cases $|f^{-1}(1)|=1$ and $|f^{-1}(2)|=1$ are disjoint. In
the first, $f^{-1}(1)=\{i\}$ with $\rho_i=\nu_1$ (and then the remaining two parts sum to
$n-\nu_1=\nu_2$ automatically); in the second, $f^{-1}(2)=\{i\}$ with $\rho_i=\nu_2$. Hence
$M_{\rho\nu}=m_{\nu_1}(\rho)+m_{\nu_2}(\rho)$.
\end{itemize}
Block inversion of a lower block-triangular matrix gives
\[
  N=\begin{pmatrix}
  D_1^{-1} & 0 & 0\\
  -D_2^{-1}A_{21}D_1^{-1} & D_2^{-1} & 0\\
  D_3^{-1}\bigl(A_{32}D_2^{-1}A_{21}-A_{31}\bigr)D_1^{-1} & -D_3^{-1}A_{32}D_2^{-1} & D_3^{-1}
  \end{pmatrix}.
\]
All entries but the $(3,1)$ block are now immediate ($D_1=(1)$, $A_{21}$ and $A_{31}$ all ones). For
the $(3,1)$ block, with $\ell(\rho)=3$,
\[
  N_{\rho,(n)}=\frac{1}{d_\rho}\Bigl(\sum_{\ell(\nu)=2}\frac{m_{\nu_1}(\rho)+m_{\nu_2}(\rho)}{d_\nu}-1\Bigr),
  \qquad d_\nu=1+\bb{\nu_1=\nu_2}.
\]
Evaluate the sum. Group the two-part partitions $\nu$ of $n$ by the unordered pair
$\{\nu_1,\nu_2\}=\{c,n-c\}$. If $c\ne n/2$ the pair contributes $m_c(\rho)+m_{n-c}(\rho)$ with
$d_\nu=1$; the single pair with $c=n/2$ contributes $2m_{n/2}(\rho)/2=m_{n/2}(\rho)$. Summing over
all pairs therefore gives $\sum_{c=1}^{n-1}m_c(\rho)$. Since $\rho$ has three parts, each at most
$n-2$, every part is counted, so the sum equals $\ell(\rho)=3$ and
$N_{\rho,(n)}=(3-1)/d_\rho=2/d_\rho$.
\end{proof}

\begin{remark}
The first column of the inverse continues in the obvious way: $1,-1/d_\rho,2/d_\rho$ for
$\ell(\rho)=1,2,3$. Theorem~\ref{thm:inverse} identifies the pattern, since the only $\pi$ with
$\lam^\pi=(n)$ is the one-block partition:
\[
  (M^{-1})_{\rho,(n)}\;=\;\frac{(-1)^{\ell(\rho)-1}\bigl(\ell(\rho)-1\bigr)!}{d_\rho}
  \;=\;[p_n]\,m_\rho .
\]
At $\rho=(1^n)$ this is $(-1)^{n-1}/n$, which is the classical coefficient $[p_n]e_n$ --- an
independent anchor on the formula at the far end of the filtration.
\end{remark}

\section{Theorem B at $L=3$, and the Bell bound}

Recall the Hall--Littlewood setting of the companion paper: $\Qp_\lam$ is the modified
Hall--Littlewood function, $Y^\lam_\rho(t)=\langle p_\rho,\Qp_\lam\rangle$, and
$\Qp_\lam=\sum_\nu c_{\lam,\nu}(t)h_\nu$. Pairing with $p_\rho$ gives
$Y^\lam_\rho=\sum_\nu M_{\rho\nu}c_{\lam,\nu}$, so every row of $M$ is a closed-form reduction of a
class of Green polynomials to $h$-coefficients. The case $\ell(\rho)=2$ is Theorem~B of that paper.
Theorem~\ref{thm:structure} and the $A_{32}$ computation give the next case for free.

\begin{theorem}[Theorem B at $L=3$]\label{thm:B3}
Let $\rho=(x,y,z)\vdash n$ with $x\ge y\ge z\ge1$. Then for every $\lam\vdash n$,
\[
  Y^{\lam}_{(x,y,z)}
  \;=\;c_{\lam,(n)}
  \;+\;\sum_{i=1}^{3}\bigl(1+\bb{2\rho_i=n}\bigr)\,c_{\lam,\ \mathrm{sort}(\rho_i,\,n-\rho_i)}
  \;+\;\Bigl(\prod_k m_k(\rho)!\Bigr)\,c_{\lam,\rho}.
\]
\end{theorem}

\begin{proof}
Expand $Y^\lam_\rho=\sum_\nu M_{\rho\nu}c_{\lam,\nu}$ and use
Theorem~\ref{thm:structure}: only $\ell(\nu)\le3$ contributes, the term $\ell(\nu)=3$ is
$d_\rho c_{\lam,\rho}$, and $\ell(\nu)=1$ contributes $c_{\lam,(n)}$. For $\ell(\nu)=2$ the
coefficient is $m_{\nu_1}(\rho)+m_{\nu_2}(\rho)$ as computed in the proof of
Theorem~\ref{thm:M3inv}. Reindexing that sum by the position $i$ rather than by $\nu$: each $i$ with
$\rho_i=\nu_1$ or $\rho_i=\nu_2$ contributes $1$, and $\nu$ is then forced to be
$\mathrm{sort}(\rho_i,n-\rho_i)$; the only overcount is at $\nu_1=\nu_2=n/2$, where a position with
$\rho_i=n/2$ is counted by both $m_{\nu_1}$ and $m_{\nu_2}$, which is the factor
$1+\bb{2\rho_i=n}$. Note $1\le\rho_i\le n-2$, so $\mathrm{sort}(\rho_i,n-\rho_i)$ really has two
positive parts.
\end{proof}

The general statement behind both theorems is about the \emph{support} of a row.

\begin{proposition}\label{prop:support}
For $\rho\vdash n$, $M_{\rho\nu}\ne0$ if and only if $\nu=\rho^\pi$ for some set partition $\pi$ of
$[\ell(\rho)]$ --- that is, if and only if $\nu$ is obtained from $\rho$ by merging parts.
Consequently
\[
  \#\{\nu:\ M_{\rho\nu}\ne0\}\ \le\ B_{\ell(\rho)},
\]
the $\ell(\rho)$-th Bell number, and a class of $\ell(\rho)=k$ parts determines $Y^\lam_\rho$ from at
most $B_k$ of the $p(n)$ coefficients $c_{\lam,\nu}$.
\end{proposition}

\begin{proof}
By Lemma~\ref{lem:threefaces} a nonzero entry is witnessed by a surjection
$f:[\ell(\rho)]\to[\ell(\nu)]$ whose fibres have $\rho$-sums $\nu_j$; its fibres form a set
partition $\pi$ with $\rho^\pi=\nu$, and conversely any such $\pi$ gives at least one $f$. The bound
is $|\Pi_k|=B_k$.
\end{proof}

\begin{example}
$B_1=1$, $B_2=2$, $B_3=5$: the one-part class sees only $c_{\lam,(n)}$, the two-part class sees two
coefficients (Theorem~B), the three-part class at most five (Theorem~\ref{thm:B3}). The bound $5$ is
attained, e.g.\ at $\rho=(3,2,1)\vdash6$, which sees
$c_{\lam,\nu}$ for $\nu\in\{(6),(5,1),(4,2),(3,3),(3,2,1)\}$ --- five of the eleven coefficients at
$n=6$. It is not attained when merges collide: $\rho=(2,2,2)\vdash6$ sees only three, because all
three two-element merges give $\nu=(4,2)$.
\end{example}

\section{Verification}\label{sec:verif}

All computations are exact: integer arithmetic for $M$, $\mathbb{Q}$ (Python
\texttt{Fraction}) for determinants and inverses, and \texttt{sympy} polynomials in $t$ for
\S6. Scripts: \texttt{scratch/2026-10-08-c1/\{mpm,verify\_structure,controls,controls2,controls3,%
inverse3,support,thmB3,ctdet\}.py}. SageMath is not installed; everything is pure Python/SymPy,
which is the point --- the symmetric-function ground truth is reached by polynomial algebra rather
than by a combinatorics library, so it is an independent mechanism and not a second reading of one
engine.

\paragraph{Falsifiable-instance census, computed before any agreement was reported.}
Claim (ii) of Theorem~\ref{thm:structure} is the new content, and it is \emph{vacuous} whenever
every length class is a singleton. That is exactly $n\le3$: the degenerate boundary of this problem.

\begin{center}
\begin{tabular}{rrrrr}
\toprule
$n$ & $p(n)$ & (i)-instances & (ii)-instances & non-vacuous (iii) ($d_\rho>1$)\\
\midrule
1 & 1  & 0   & \textbf{0}   & 0\\
2 & 2  & 1   & \textbf{0}   & 1\\
3 & 3  & 3   & \textbf{0}   & 1\\
4 & 5  & 9   & 2   & 3\\
5 & 7  & 19  & 4   & 4\\
6 & 11 & 48  & 14  & 7\\
7 & 15 & 92  & 26  & 10\\
8 & 22 & 201 & 60  & 16\\
9 & 30 & 379 & 112 & 22\\
\bottomrule
\end{tabular}
\end{center}

\paragraph{A. Two engines for the entry.} Engine~A counts the maps of
Lemma~\ref{lem:threefaces} (dynamic programming over the parts, itself checked against brute-force
enumeration of all $\ell(\nu)^{\ell(\rho)}$ maps: $209$ pairs, $0$ mismatches). Engine~B multiplies
out $p_\rho=\prod_i(\sum_v x_v^{\rho_i})$ as a polynomial and reads the coefficient of
$x_1^{\nu_1}\cdots x_k^{\nu_k}$. Over all $n\le9$: \textbf{1818 pairs compared, 736 nonzero, 0
mismatches}.

\paragraph{B. The structure theorem.} Direct check of (i)/(ii)/(iii) for all $n\le9$: $752$
instances of (i), $218$ instances of (ii), $64$ non-vacuous instances of (iii) --- \textbf{0
failures}.

\paragraph{C. The determinant.} Corollary~\ref{cor:det} against exact Gaussian elimination, for
\emph{every} $(n,L)$ with $n\le9$: \textbf{0 mismatches}. $L=1$ returns $1$ in every case, the free
sanity check. Corollary~\ref{cor:unimod}'s characterisation of $\det=1$: all $(n,L)$ with $n\le10$,
$0$ mismatches.

\paragraph{D. The determinant, through a different matrix.} $|\det(\text{character table of }S_n)|$
computed by Murnaghan--Nakayama (validated by column orthogonality
$\sum_\rho(\chi^\lam_\rho)^2/z_\rho=1$ for every $\lam$) equals $\prod_{\rho\vdash n}d_\rho$ for
$n\le7$: $1,2,6,96,2880,9953280,100329062400$.

\paragraph{E. The inverse.} Theorem~\ref{thm:inverse} against Gaussian elimination: $434$ entries,
$n\le7$, $204$ nonzero, $0$ mismatches. Corollary~\ref{cor:truncinv}: $(M^{(L)})^{-1}$ versus the
truncation of $M^{-1}$, all $(n,L)$ with $n\le7$, $0$ mismatches. Corollary~\ref{cor:denom}: for
every row $\nu$ ($44$ rows, $n\le7$) the l.c.m.\ of the denominators is exactly $d_\nu$.
Theorem~\ref{thm:M3inv} against Gaussian elimination: \textbf{608 entries, $n\le10$, 192 nonzero, 0
mismatches}.

\paragraph{F. Theorem~\ref{thm:B3}, through a chain that never counts a set partition.}
$\Qp_\lam=\sum_\nu K_{\nu\lam}(t)s_\nu$ with $K$ from the \emph{charge statistic} on semistandard
tableaux; $Y^\lam_\rho=\sum_\nu K_{\nu\lam}(t)\chi^\nu_\rho$ with $\chi$ from
Murnaghan--Nakayama; $c_{\lam,\nu}=\sum_\mu K_{\mu\lam}(t)\,[h_\nu]s_\mu$ with $[h_\nu]s_\mu$ from
the Jacobi--Trudi determinant (validated by $M(s,h)M(h,s)=I$ for $n\le7$). No step uses
$M_{\rho\nu}$. Result: \textbf{115 checks, $3\le n\le7$, all with $Y^\lam_\rho\ne0$, 0 failures}.
Proposition~\ref{prop:support}: support $=$ coarsenings, all $66$ rows with $n\le8$, $0$ mismatches.

\paragraph{G. Controls.} Each claim gets a planted control, and each silence is diagnosed.
\begin{center}
\small
\begin{tabular}{llp{5.4cm}}
\toprule
control & reading & note\\
\midrule
C0 drop ``fibres nonempty'' & \textbf{silent}, $0/434$ & \emph{Vacuous, and provably so}: by
Remark~\ref{rem:nonempty} the guard is implied by $\nu_j\ge1$, so flipping it cannot change any
entry. Replaced by C1.\\
C1 force blocks to be intervals & \textbf{fires} $139/434$ & smallest witness
$\rho=\nu=(1,1)$: $1$ vs $2$.\\
C2 $d_\rho\rightsquigarrow z_\rho$ in Cor.~\ref{cor:det} & \textbf{fires} $27/28$ & first at
$(n,L)=(2,1)$.\\
C3 $m_i!\rightsquigarrow m_i$ & \textbf{fires} $17/36$ & could only fire where some $m_i\ge3$;
census says $17$ of $36$ slots, and it fires at all $17$. First at $(3,3)$.\\
C4a plant one nonzero in a same-length off-diagonal slot & checker \textbf{fires} ($1$ failure);
$\det$ \textbf{silent} ($24\to24$) & \emph{Vacuous for the determinant, and provably so}: a single
off-diagonal perturbation of a diagonal block leaves it triangular, so $\det$ is unchanged. The
consequence matters: \textbf{the determinant agreement in C is evidence for (iii) only, conditional
on (i)+(ii)} --- it is not independent evidence for (ii).\\
C4b plant a symmetric pair & \textbf{fires} $\det 24\to12$ & the control the determinant can see.\\
C5a Möbius sign flipped & \textbf{fires} $95/434$ &\\
C5b $(|B|-1)!$ dropped & \textbf{fires} $83/434$ & first at $\lam=(1,1,1)$, where a block of size
$3$ first exists.\\
C6 transpose, truncation vs inversion & \textbf{silent} & \emph{Vacuous, and provably so}: leading
principal truncation commutes with inversion for \emph{every} triangular matrix, upper or lower, and
$M^{\mathsf T}$ is upper triangular. Replaced by C6$'$/C6$''$.\\
C6$'$ plant one cross-length entry above the length diagonal & \textbf{fires} & two of three plants
make the full matrix \emph{singular} while all $M^{(L)}$ stay invertible, so ``all leading blocks
invertible'' needs triangularity; the third gives a genuine truncation/inversion mismatch.\\
C6$''$ character table, truncate by $\ell(\lam)$ & \textbf{fires} $243/368$ at $n=6$ & the
commutation test is sensitive; it is not a tautology of the harness.\\
C7 admit zero parts in $\nu$ & \textbf{fires} & Remark~\ref{rem:loadbearing}: a fact about the
statement.\\
C8 drop $(1+\bb{\nu_1=\nu_2})$ in Thm.~\ref{thm:M3inv} & \textbf{fires} $6/116$ & census says
exactly $6$ slots could fire ($n$ even, $m_{n/2}(\rho)>0$); it fires at all $6$. First at
$n=4$, $\rho=(2,1,1)$, $\nu=(2,2)$.\\
\bottomrule
\end{tabular}
\end{center}

\paragraph{An instrument fault, recorded because it read as a refutation.}
The first run of F reported \textbf{86 of 115 failures}. The polynomials were equal; the comparison
was not. In the Murnaghan--Nakayama recursion the sign was coded as \verb|(-1)**ht| with $ht$ the
signed displacement of a $\beta$-number, which is \emph{negative}; in Python \verb|(-1)**(-1)| is
the float $-1.0$, so $\chi^\nu_\rho$ was returned as a float, $Y$ carried floats, and
\verb|sympy|'s \verb|!=| separated $t^3+t^2+t-1.0$ from $t^3+t^2+t-1$. Every value was correct and
every comparison was wrong. Had I read that output as data, I would have recorded a true theorem as
refuted at $n=4$. The repair is \verb|1 if ht % 2 == 0 else -1|, after which the character engine
passes column orthogonality at $n=5,6$ and F is clean. \emph{A disagreement between two
representations of one value is not a disagreement about the value} --- and a $0$ failure count is
only meaningful once the comparison operator has been shown to compare what it claims to.

\section{What is and is not new}

The $m$-to-$p$ transition (Theorem~\ref{thm:inverse}) is classical, as is the evaluation of
$\langle p_\rho,h_\nu\rangle$ as an ordered-set-partition count (Lemma~\ref{lem:threefaces}, which
is Lemma~2.3 of the companion paper in the two-part case). The determinant of the $S_n$ character
table is also classical, and by the remark after Corollary~\ref{cor:unimod} it is equivalent to the
$L=n$ case of Corollary~\ref{cor:det}. I have not searched the literature this session (no browsing),
so I claim no novelty for any of those three.

What this session adds is the \emph{filtered} picture and its consequences:
Theorem~\ref{thm:structure}(ii) --- the diagonal blocks are diagonal, not merely invertible ---
and with it the closed form for $\det M^{(L)}$ at every truncation
(Corollary~\ref{cor:det}), the exact unimodularity criterion (Corollary~\ref{cor:unimod}), the exact
denominators of the duality (Corollary~\ref{cor:denom}), the explicit $(M^{(3)})^{-1}$
(Theorem~\ref{thm:M3inv}), the three-part analogue of Theorem~B (Theorem~\ref{thm:B3}), and the
Bell-number bound on how many $h$-coefficients a $k$-part class can see
(Proposition~\ref{prop:support}).

\paragraph{Gaps.} None in the proofs above; every step is an identity, a definition, or Möbius
inversion in $\Pi_k$. Three things are \emph{not} done here and should not be read as done:
\begin{enumerate}
\item No closed form for $c_{\lam,\nu}$ at $\ell(\nu)=3$ is proved. Theorem~\ref{thm:B3} reduces the
three-part Green polynomial to five $h$-coefficients; it does not evaluate them. The missing input is
a three-point analogue of Rick's two-point shuffle formula, and whether $Sh_{A,B}$ composes to a
triple shuffle is untouched by this session.
\item Nothing here is formalised in Lean.
\item The claim that Lemma~\ref{lem:threefaces} and Theorem~\ref{thm:inverse} are classical is a
recollection, not a literature check; no browsing was done.
\end{enumerate}

\end{document}
